\section{Introduction}\label{sec:intro}
%LinkedIn, as the world's largest professional network, serves over a billion members globally, providing services ranging from job searches to content engagement. This paper offers a comprehensive view on {\systemname}, LinkedIn's expansive model-based GPU retrieval system. We delve into the construction of high-quality search and recommender systems, focusing on embedding-based retrieval (EBR). 

%Traditional EBR relies on unsupervised nearest neighbor search solutions, indexing item vectors for faster retrieval \cite{hnsw_paper, pq_paper}. Our paper introduces an innovative approach of integrating exhaustive search with pre-filtering into a differentiable model on GPU, using neural networks for distance learning and ranking. In {\systemname}, item vectors and model weights coexist within the same model binary, diverging from traditional search indexing methods. Our perspective is that the future of search and recommender systems will converge in differentiable model-based serving. This approach allows for the joint optimization of retrieval and ranking systems, merging the boundaries between these systems. 

%The essential embedding-based retrieval method, the K-nearest neighbor (KNN) search algorithm, uses learned query and item embeddings with a specific similarity metric. It scans the item set, calculates the similarity score for each item relative to the query, and selects the top-K closest items as candidates. Typically, KNN uses dot-product similarity, a form of matrix multiplication with normalized embeddings, which has been significantly sped up on modern GPUs (A100, H100, etc.) in frameworks like PyTorch and TensorFlow. Several challenges motivate us to propose model-based KNN algorithms implemented on GPUs:

LinkedIn, the world's largest professional network, serves over a billion members globally, offering services from job searches to content engagement. This paper explores {\systemname}, LinkedIn's model-based GPU retrieval system, focusing on embedding-based retrieval (EBR). Traditional EBR uses unsupervised nearest neighbor search solutions~\cite{hnsw_paper, pq_paper}, indexing item vectors for fast retrieval. Our paper presents an innovative approach, combining exhaustive search with pre-filtering in a differentiable GPU model, using neural networks for distance learning and ranking. In {\systemname}, item vectors and model weights coexist within the same model binary, unlike traditional search indexing methods.

We believe the future of search and recommender systems lies in differentiable model-based serving, enabling joint optimization of retrieval and ranking. The K-nearest neighbor (KNN) search algorithm, an essential embedding-based retrieval method, uses learned query and item embeddings with a specific similarity metric to select the top-K closest items. Typically, KNN uses dot-product similarity, a form of matrix multiplication with normalized embeddings, which has been significantly sped up on modern GPUs (A100, H100, etc.) in frameworks like PyTorch and TensorFlow. Several challenges motivate us to propose model-based KNN algorithms implemented on GPUs:
%\vspace{-4pt}
\begin{itemize}[leftmargin=*]
    \item Liquidity challenge: Real-time search systems rely on specific attributes to filter relevant items. In job recommendation systems, for example, filters like company names, locations, and skills are essential. Items meeting these conditions must be prioritized to avoid exclusion due to low KNN scores from embeddings alone.
    %\item Liquidity challenge: Real-time search systems depend on definite attributes to filter relevant items per query. For instance, in job recommendation systems, postings include company names, locations, and skill requirements—explicit filters users might specify. Items meeting these query conditions must be prioritized in rankings to prevent exclusion due to low KNN scores from embeddings alone.

    \item Low latency requirement: Reducing retrieval latency and increasing throughput is a constant priority.
    %Low latency requirement: reducing the retrieval latency and increasing the throughput is an enduring topic.

    \item Huge memory cost: As number item embeddings and clauses increase, finding ways to lower memory usage and boost computational speed without compromising retrieval quality presents a significant challenge.

    \item Freshness: Demonstrate that model-based approaches can enhance traditional nearest neighbor searches in quality and latency while supporting functionalities like live updates.
\end{itemize}


In this paper we discuss deployment of large-scale, neural model-based retrieval system, highlighting key challenges and solutions. A major challenge was the absence of efficient pre-filtering in PyTorch and TensorFlow, addressed by our custom indexing and filtering methods detailed in \S\ref{sec:knn}, which also tackle latency issues. We also cover memory cost management through quantization techniques for larger indexes in \S\ref{sec:quant_knn}. {\systemname} enhanced search quality, utilizing multi-embedding retrieval algorithms discussed in \S\ref{sec:mexture_similarities}. Our work positions us among the pioneers in the industry in introducing a retrieval model-based serving infrastructure (\S\ref{sec:MLinfra}), showcasing the capability of such model-based retrieval systems to be effectively live-updated at scale (\S\ref{sec:live_update}). We perform our study of model-based index serving focuses on interest-based recommendations on LinkedIn's Feed, also known as out-of-network (OON) recommendations. These recommendations leverage member profiles and previous interactions with the Feed, enabling LinkedIn members to access highly relevant content. We integrate OON content into various LinkedIn surfaces, like Feed and Notifications, based on predicted user engagement likelihood. The effectiveness of OON recommendations is gauged by member interactions with OON content. We use two-tower neural networks to create embeddings for members and Feed Posts, forming a candidate selection vertical for OON in the Feed through EBR with a differentiable model-based search index. {\systemname} significantly outperforms FAISS-based~\cite{FAISS_paper} retrieval system in OON recommendations. We support full-scan model-based index serving on GPUs with latencies as low as 4 ms, handling indexes from 15 million to a billion entries. This capability, along with modeling enhancements, significantly boosts quality as detailed in \S\ref{sec:A_b_test}.